% 本文件由 LaTeX 排版 Agent 根据有效 translation.json 自动生成。
% author=Gregory of Nyssa; work_id=2904; title=论天主圣三

\chapter{论天主圣三}

% structure: p0001 source=h1 target=omitted reason=page-title-already-rendered-by-parent

致\Person{欧斯塔削}。\par

所有研习医学的人，可以这么说，都是以仁爱为职业；我想，若有人将你们的学问置于人生一切严肃追求之上，他必作正确的判断，不致失于明断，因为生命——人最珍视的一切——若是没有健康，便成为可逃避的、充满痛苦的事，而健康正是你们的技艺所供给的。但在你们身上，这学问具有显著的双重功效：你们为其仁爱拓宽了界限，因为你们不仅将技艺的益处限于人的身体，也关注心灵疾病的医治。我这样说，不仅因为众口相传，而且因我亲身体验，在许多事上，尤其在这次我们的敌人那难以形容的恶意中，你们巧妙地驱散了它，它像一股邪恶的洪水席卷了我们的生命，你们用安慰的话语消解了我们心中剧烈的炎症。我原以为，面对敌人持续不断的各种攻击，应当保持沉默，安静地承受他们的袭击，而不是与那些以谎言为武器的恶人对质，那武器最有害，有时甚至将矛头刺向真理。但你们劝我不要背叛真理，而要驳斥诽谤者，以免谎言胜过真理，使多人受害，这实在是善举。\par

我可以说，那些无端仇视我们的人，其行为与\WorkTitle{伊索寓言}中的原则颇为相似。正如他让狼对羔羊提出一些指控（我想，狼也羞于无故毁灭一个不曾伤害它的无辜者），然后，当羔羊轻易驳倒了所有诽谤的指控时，狼却毫不放松攻击，反以牙齿取胜，尽管公理已败；同样，那些像追求美善一样热切仇恨我们的人（或许也羞于显得无故仇恨），便捏造对我们的指控和抱怨，但他们并不坚持自己所说的任何话，而是时而说此事，片刻后又说彼事，然后又以另一事作为他们敌视我们的原因。他们的恶意不立足于任何根据，一旦从一项指控中被逐出，便紧抓另一项，再从那里又抓住第三项，即使所有指控都被驳倒，他们仍不放弃仇恨。他们指控我们宣讲三位天主，并向众人耳边灌输这诽谤，从不停止，千方百计让人相信。这时，真理与我们并肩作战，因为我们公开向所有人，私下与那些同我们交谈的人表明，我们咒骂任何说有三位天主的人，并视之为非基督徒。然后，他们一听到这话，便拿\Person{撒伯流}作攻击我们的方便武器，他散布的瘟疫成了持续攻击我们的主题。我们又一次以惯用的真理武装来抵抗这攻击，表明我们同样憎恶这种异端形式，正如憎恨犹太教一样。那么，他们经过这番努力后是否疲倦，安于休息？绝非如此。如今他们指责我们创新，并以这种方式提出对我们的控诉：——他们声称，当我们承认三位时，却说有一位美善、一位能力、一位天主性。这种说法并不越出真理，因为我们确实这样说。但他们抱怨的理由是，他们的习惯不接纳此说，圣经也不支持。那么我们如何回答？我们认为，不应以他们的通行习惯作为纯正教义的法律和准则。因为如果习惯可以证明教义的正确，那么我们当然也可以提出我们的通行习惯；而如果他们拒绝我们的习惯，我们当然也不必遵从他们的习惯。那么，让默感的圣经作我们的仲裁者吧，真理的裁决必然归于那些教义与圣言相符的人。\par

那么，他们的指控是什么？在对我们的控告中，同时提出了两点：一是我们分开各位格；二是我们不使用任何复数形式的归属于天主的名词，而（如我所说）以单数形式谈论美善、能力、天主性以及所有这类属性。关于分开各位格，那些主张在神圣本性中有不同本体的人，不能很好地反对。因为那些说有三个本体的人，不可能不说有三位格。所以，只有这一点受到质疑：那些归属于神圣本性的属性，我们以单数形式使用。\par

但我们对此的回应是现成的、明确的。因为任何谴责我们说天主性是一的人，必然要么支持那些说有很多个的人，要么支持那些说没有的人。但是默感的教导不容许我们说有很多个，因为每当它使用这个术语时，都以单数形式提及天主性，如：\BibleQuote{在祂内住有整个天主性体的圆满}\BibleRef{哥 2:9}；以及另一处：\BibleQuote{祂那看不见的事，从创世以来，藉着所造之物，就可明白察知，即祂永恒的能力和天主性}\BibleRef{罗 1:20}。因此，如果把天主性的数目扩大到多数，仅属于那些患有多神谬误之疫的人，而彻底否认天主性则是无神论者的教义，那么，指责我们说天主性是一，是什么教义呢？但他们更清楚地揭示了他们论证的目的。关于圣父，他们承认祂是天主，圣子也享有天主性的尊荣；但圣神，与圣父和圣子同列，他们却不能将其纳入天主性的概念中，而认为天主性的能力从圣父发出到圣子，在那里停止，从而将圣神的性情与神圣荣光分开。因此，我们必须在尽可能短的篇幅内，也回应这个意见。\par

那么，我们的教义是什么？主在将救恩的信德交付给那些成为圣言的门徒时，也将圣神与圣父及圣子联合在一起；我们确认，一旦联合，便是持续的联合；因为它并非在某一事上联合，而在其他事上分离。但圣神的德能，在赋予生命的德能中——藉此德能我们的本性从可朽坏的生命转移到不朽——并与圣父及圣子一同被包含在内，而且在许多其他情形中，如在\Quote{善}、\Quote{圣}、\Quote{永恒}、\Quote{智慧}、\Quote{正义}、\Quote{首领}、\Quote{大能}的观念中，实际上处处如此，它在一切以特殊卓越意义归于天主的属性上，与它们有不可分离的关联。因此我们认为，正确的想法是：那在如此崇高卓越的观念中与圣父和圣子联合的，在任何方面都不应与它们分离。因为我们不知道，在表达我们对天主性体之概念的属性中有任何高低之分，致使我们应当认为（虽允许圣神在较低属性上有共融）判断祂不配那些较高属性是一种虔诚的行为。因为一切天主属性，无论是被命名或构想，彼此都是同等地位的，因为它们在其主体之意指方面是不可区分的。因为\Quote{善}这个称呼并不引导我们的心智到一个主体，而\Quote{智慧}、\Quote{大能}或\Quote{正义}的称呼到另一个主体，但所有这些属性所指的对象是同一个；并且，若你谈论天主，你所指的就是你由其他属性所领悟的那同一位。如果所有归于天主性体的属性在指示主体方面力量相等，从不同方面引导我们的心智到同一主体，那么有什么理由让人一方面允许圣神在其它属性上与圣父和圣子共融，另一方面却唯独将祂排除在天主性之外？绝对必须要么也允许祂在此属性上的共融，要么不承认祂在其它属性上的共融。因为如果祂在那些属性上是有资格的，那么祂在此属性上肯定也不较少有资格。但若祂按他们的说法是\Quote{较小}，以致于祂被排除在与圣父和圣子在\Quote{天主}属性上的共融之外，那么祂同样不配分享任何其他属于天主的属性。因为那些属性，当被正确领悟并藉我们在每种情形中默观的概念相互比较时，将被发现所意指的并不亚于\Quote{天主}这个称呼。而这一点的一个证明是，甚至许多低等存有也被称为这个名字。此外，圣经并不吝于将这个名称用于不相称的事物，例如它用\Quote{天主}这个称呼来指称偶像。因为经上说：\BibleQuote{愿那没有创造天地的众神消灭，被投到地底下}；以及\BibleQuote{外邦人的众神都是魔鬼}；而那个巫婆在咒语中，当她为撒乌耳召来他所寻求的亡魂时，她说她\BibleQuote{看见了神}（\BibleRef{撒上 28:13}）。又比如，巴郎是一个卜者兼先见者，从事占卜，为自己获得了魔鬼的教导和巫术占卜，圣经说他从天主那里得到指示。\EditorialNote{户籍纪第二十二章}人们可以通过从圣经中收集许多同样的事例来证明，这个属性并不优越于其他属于天主的属性，因为正如所说，我们发现它以一种歧义的方式被用于不相称的事物；但我们在圣经中从未被教导说，\Quote{圣}、\Quote{不朽}、\Quote{正义}、\Quote{善}这些名称会被用于不相配的事物。那么，如果他们不否认圣神在那些以特殊卓越意义虔诚地仅归于天主性体的属性上与圣父和圣子有共融，有什么理由假装祂唯独在这一属性上被排除，而在此属性上，已表明连魔鬼和偶像也以一种歧义的方式分享呢？\par

但他们说，这个称谓指明性体，并且由于圣神的性体与圣父和圣子不共通，因此祂也不分享这个属性的共通性。那么，请他们指出来，他们凭借什么识别出这种性体的差异。因为如果有可能在绝对本质中默观天主性体，并且我们能够凭显现象看出什么属于它、什么不属于它，那么我们就确实不需要其他论据或证据来理解这个问题。但由于它超越提问者的理解，我们不得不从某些特定证据出发，就那些超出我们认知的事物进行论证，因此，我们绝对有必要通过天主的行动来引导我们探讨天主性体。那么，如果我们看到由圣父、圣子和圣神所实行的行动彼此不同，我们就会从行动的不同特性推测出行动者的性体也不同。因为那些在本性上不同的事物，不可能在行动方式上一致：火不会冷却，冰不会供热，而是它们的行动随着本性差异而区别。另一方面，如果我们理解圣父、圣子和圣神的行动是一个，没有任何不同或变化，那么就必须从行动的同一种推知它们性体的同一。圣父、圣子和圣神同样赐予圣化、生命、光明、安慰以及一切类似的恩宠。并且，当有人在福音中听到救主对圣父论及祂的门徒时，不要将圣化的能力特别归于圣神，祂说：\Quote{父啊，请因祢的名圣化他们。}同样，所有其他恩赐也由圣父、圣子和圣神在那些堪当的人身上施行：一切恩宠和能力、引导、生命、安慰、转变为不朽、通往自由，以及降临到我们身上的每一种其他恩惠。\par

但是，那在我们之上的事物秩序，无论在那理智领域还是在那感官领域（如果我们能借着我们所知道的，对那些在我们之上的事物也作出推测），它本身是在圣神的行动与大能内建立的，每个人都按自己的功德与需要领受恩惠。因为虽然那在我们本性之上的事物的安排与秩序，对我们的感官是隐晦的，然而我们借我们所知的事物，更合理地推断圣神的大能也在它们之中工作，而不认为它被排除在我们之上事物的现存秩序之外。因为那主张后一种观点的人，是以赤裸裸、不体面的方式提出他的亵渎之语，无法用任何论据支持他那荒谬的意见。但同意我们之上的事物，也由圣神与圣父和圣子一同以祂的大能来安排的人，在这一点上，是从他自己的生命中得到了明确的证据来支持他的论断。因为如同人的本性是由身体与灵魂所构成，而天使的本性则分有非身体的生命，如果圣神只在身体的情况中工作，而灵魂不能接受那来自祂的恩宠，那么，或许可以由此推断：如果我们内那理智的与非物质的本性是在圣神的大能之上，那么天使的生命也同样不需要祂的恩宠。但如果圣神的恩赐主要是灵魂的恩宠，而灵魂的构成藉其理智性与不可见性而与天使的生命相连，那么，一个知道如何看出推论的人，怎能不同意：每一个理智的本性都是受圣神的安排所治理的呢？因为经上说：\BibleQuote{天使在天上，常看见我父的面}\BibleRef{玛 18:10}，而不藉着祂的肖像凝视祂，就不可能看见圣父的位格；而圣父位格的肖像，就是那独生子；凡心智没有被圣神光照的人，也不能接近祂；由此还显示了什么呢？岂不是圣神没有与圣父和圣子所行的任何行动分离吗？因此，圣父、圣子、圣神在行动上的同一性，清楚地显示了它们本体的不可分辨的特性。所以，即使\Quote{天主性}这个名称确实指明本性，本体的共通性也显示，这个称号也同样正当地适用于圣神。但我不晓得这些制造各种论据的人，如何引用\Quote{天主性}这个称号来指明本性，好像他们没有从圣经中听到，这是一个\OriginalTerm{任命}{appointment}的事情，而本性并非由此而生。因为梅瑟被任命为埃及人的神，因为那赐给他神谕者等，这样对他说：\BibleQuote{我已立你为法郎的神}\BibleRef{出 7:1}。因此，这个称号的力量是表示某种权能，或是监督的权能，或是行动的权能。但天主本性本身，其真实所是，是超越所有为它设想的名称的，正如我们的教义所宣示的。因为当得知祂是慈善的、审判者、良善的、公义的，以及其他一切同类属性时，我们所学到的是祂行动的多样性，但凭我们对祂行动的认识，我们并不能更多地认识那行动者的本性。因为当一个人对这些属性中的任何一个，以及对应用这些名称的本性给出定义时，他不会对两者给出相同的定义：而定义不同的事物，其本性也是不同的。确实，本体是一回事，没有任何定义能够表达它；而有关它所用的名称的意义则随情况而变化，因为这些名称是从某种行动或属性而起的。如今，我们从属性的共通性得知行动中没有区别，但对于本性方面的差异，我们找不到明确的证据，因为行动的同一性，正如我们所说，反而指明了本性的共通性。那么，如果\Quote{天主性}是一个源自行动的名称，正如我们说圣父、圣子、圣神的行动是同一的，所以我们也说天主性是同一的：或者，如果按照多数人的看法，天主性是指明本性的，既然我们在它们的本性中找不到任何不同，我们不无理由地将天主圣三定义为同一个天主性。\par

但若有人称此名号指示尊位，我不明其凭何推理将此词引申至此义。然而，因常闻多人言及此类说法，为使反对者的热忱无从攻击真理，我们姑且随持此见者偏离正途，考量此种意见，并言：即使此名确指尊位，在此情形下，此名号亦将恰当地适于圣神。因为王权的属性指向一切尊位；经上却说：\Quote{我们的天主}是\Quote{从永远就为君王}。然而圣子拥有属于圣父的一切，自身亦被圣经宣告为君王。而神圣圣经称圣神为独生者的傅油（\BibleRef{宗 10:38}），藉借用此世常用之词以解释圣神的尊位。因为古时，在那些被提升至王权的人身上，此尊位的标记是施于他们的傅油；当此事发生后，他们便从私人与卑微的地位转变为统治的优越，而被视为配得上此恩宠者，在受傅后另得一名，不再称常人为\Quote{主的受傅者}。为此缘故，为使圣神的尊位更明确地向人显示，圣经称祂为\Quote{王国的标记}和\Quote{傅油}，由此我们得知，圣神共享天主独生圣子的荣耀与王国。因为在以色列，未先受傅者不得进入王权；同样，圣经藉借用我们常用之词，表明权能的平等，显示甚至圣子的王国也未经圣神的尊位而领受。为此缘故，祂被合宜地称为基督，因为此名给出祂与圣神不可分、不可离之结合的明证。如此，若独生天主是受傅者，圣神是祂的傅油，且受傅者之名指向王权，傅油是其王权的标记，那么圣神亦共享祂的尊位。因此，若他们说天主的属性指示尊位，而圣神被显示共享此属性，则分享尊位者也必分享代表尊位之名。\par

\SourceRecord{Translated by H.A. Wilson.；Nicene and Post-Nicene Fathers, Second Series；Vol. 5.；Edited by Philip Schaff and Henry Wace.；Buffalo, NY: Christian Literature Publishing Co.,；1893.；<http://www.newadvent.org/fathers/2904.htm>.}
